\section{A Finite FPGA Mapping and Its Validation}
\subsection{Connected Architecture}
The FPGA implements ALT sample execution, not STRICT's native flood order. Its complete core connects 64 candidate lanes in 16 frontend groups, 64 canonical banks and a shared 32-port memory network. Samples execute sequentially at the batch controller. Runtime configuration supplies source identities and eligible-target counts, never expected answers. The admitted core domain includes at most $2^{22}$ nodes, 32,768 cells and 256 samples, with checked 32-bit line addresses and disjoint, aligned regions. Host admission validates descriptor/edge encodings and cost indices; this is an explicit invariant, not a second hardware check inside every leaf.

The controller performs real per-sample state initialization, cost loading and source seeding. A finite directory/minimum tree selects work; token reservation precedes canonical POP. The frontend keeps transaction, descriptor, ready and target state finite, with registered request/return contexts and backpressure. Exact packed descriptors supply target identity and edge ranges. Four lanes per group form FP32 whole-pair proposals. A connectivity-specialized network routes only legal request/return connections while preserving source, kind, generation and local identity; finite credits and tags govern physical returns.

Each bank reserves an in-flight set through proposal read, comparison and writeback. Pipelined occupied-set enumeration feeds a finite scratch path: unspilled resident fragments can bypass a scratch read, whereas spilled fragments perform the required merge. Canonical storage supports eight live transactions per bank with same-word exclusion and complete-word updates. Directory publication remains serialized through actual response and deferred acknowledgment; accepting a fragment does not retire its canonical effects. Expansion fences wait for frontend, candidate, bank and network drain. Finalization, directory reopening and output writes are charged, and completion waits for the final write acknowledgments. These finite services constrain throughput even when candidate arithmetic can accept work every cycle.

The final RTL revision changes only return-payload selection to a 116-bit bank-selected word multiplexer. The other 30 DUT source files and sequential state logic are unchanged. It adds no lane, queue capacity or modeled service advantage. This isolated revision tests whether reducing selection logic improves physical feasibility; reduced logic alone is not a timing result.

\subsection{RTL and Model Evidence}
The complete production-size RTL executes two native graphs with all 8 and all 12 samples. Under the baseline test service, they take 157,312 and 268,916 cycles, with 42 and 106 waves and 4,192 and 21,511 physical requests, respectively. Every request is acknowledged and complete outputs equal the frozen ALT references. A full second-graph stress run takes 276,758 cycles with the same outputs and transaction counts. Stress adds deterministic 0--16-cycle return jitter and admits requests during 14 of every 19 cycles. Eleven invalid configurations are rejected without a physical request.

Across baseline, stress and conservative service, calibration covers 52 complete executions. Total cycles, all 19 mutually exclusive phases, 3,772 ordered intervals, 402 waves and 72,917 request/return events match the finite model, including addresses and identities. A prior mismatch was repaired by modeling a real registered frontend-idle visibility stage, not by fitting a cycle offset. Local protocol, selector, rejection and lifecycle checks support the revised multiplexer. These are executed-case checks, not universal formal equivalence or post-synthesis equivalence proofs.

The full BNN40/LSTM40 evaluation uses the bound finite model: 160 complete sample outputs across two workloads and two service scenarios, with exact output hashes, phase sums, request/return conservation and finite endpoint capacity checks. Reuse is justified by the unchanged source/service binding and the new RTL calibration. Small-graph cycles are not scaled by graph size or lane count to obtain full-batch performance.
